\section{Threats to Validity}\label{sec:limitations}

\subsection{Construct validity}

The seven antipattern classes follow study-specific operational definitions derived from Karwin's catalogue \nobreakcite{karwin2023sql}. Different boundaries for sentinel values, implicit projections, search patterns, or inferred keys would change the reference and detector counts. Agreement with the single-annotator reference therefore measures conformance to these definitions and one annotator's application of them. Annotation correctness against an independent consensus remains unknown.

The primary matching rule treats class-labelled line spans as events and requires one-to-one overlap at intersection over union \num{0.50}. This construct captures class, occurrence count, and approximate localisation. It does not measure semantic equivalence between differently grouped occurrences, and true negatives are undefined. The small threshold sensitivity supports robustness to the tested span boundary alone.

The corpus measures have narrower meanings than prevalence or risk. They count detector flags and repositories with at least one flag without correcting for class-specific error, repository size, or relevant API exposure. The post hoc categories classify returned source fragments without resolving call targets. They therefore describe recurring forms among flags, while prevalence and API-associated risk remain unmeasured. The study evaluates algorithmic agreement. Developer comprehension, trust, action, and practical utility remain unmeasured.

\subsection{Internal validity}

The primary evaluation retains the original annotations, but the sensitivity analysis revised records after inspecting detector disagreements. This detector-informed process can preferentially recover errors visible to the detector while leaving shared omissions undiscovered. The revised-reference analysis is consequently optimistic.

Class retention and project partitioning were data-dependent. The support rule used the complete \num{61}-repository annotation set, and the split was selected from \num{1000000} seeds using per-class occurrence balance. Project-disjoint assignment prevents direct file overlap among training, validation, and test sets, yet the retained classes and split composition remain adapted to the observed sample.

The preliminary sampling frame also contained one duplicate repository and used an incomplete path list for excluding irrelevant generated files. Later corrections improved the corpus run but could not change the already drawn sample and partitions. In addition, model and prompt selection used one validation partition and one observed run per configuration. The selected configuration may therefore reflect idiosyncrasies of that partition or execution.

\subsection{External validity}

The target population is the set identified by a dated GitHub search and the applied filters. The search can omit repositories outside GitHub, projects without Maven or Gradle declarations, repositories omitted when a query exceeds the API's \num{1000}-result cap, and projects that do not retain generated jOOQ classes. Claims from the \num{602}-repository corpus apply to this identified population. Coverage of all jOOQ software is unknown.

All analysed repositories are open source. Proprietary systems may differ in schema scale, coding conventions, review practice, and application domain. The detector evaluation also covers seven supported classes from the original 19-class codebook and only Java code using jOOQ. Transfer to excluded antipattern classes, other query builders, and runtime-only manifestations is unknown.

The \num{602}-repository corpus includes all \num{61} sampled repositories and is therefore not an independent transfer population. Repositories outside that sample have no reference annotations, and the corpus as a whole has no independent audit. Held-out agreement may not transfer to projects with source forms absent from the reference sample. A transfer audit that includes files with and without flags is required before interpreting the corpus output as calibrated detection performance.

\subsection{Conclusion validity}

The held-out reference is imbalanced across classes and projects. Implicit Columns contributes \num{270} of the \num{523} reference occurrences, so the micro aggregate is dominated by that class. Project support is especially sparse for 31 Flavors in one test repository, Rounding Errors in two, Keyless Entry and Poor Man's Search Engine in three each, and Fear of the Unknown in five. Per-class point estimates for these classes are therefore sensitive to individual repositories.

The project-cluster bootstrap preserves within-project dependence, but its percentile range is conditional on the label-informed \num{20}-project test partition. It is neither a population confidence interval nor an estimate of run-to-run variation. Class-level bootstrap ranges were unsuitable because many resamples omitted all support for sparse classes.

Configuration comparisons are descriptive results from one run per model-prompt pair on one validation partition. They support the reported selection decision without establishing statistically reliable rankings. Likewise, the intersection-over-union sensitivity analysis varies only the span threshold; it does not test annotation correctness, alternative matching constructs, or model stochasticity. Conclusions about class differences and configuration performance should remain bounded to the observed runs.

\subsection{Reliability validity}

The repeat-annotation exercise measures temporal consistency by the same annotator after a 30-day interval. Its F1-score of \num{0.925} does not provide independent reproduction or adjudicated correctness. A second annotator applying the frozen codebook would be needed to assess whether the decisions can be reproduced across people.

The selected detector was run once on the held-out partition and once on the corpus. Its run-to-run variance is therefore unknown. The hosted commercial model and gateway can also change over time, so a stable model name and fixed parameters do not guarantee identical future outputs. Reproduction depends on continued service access as well as preserved local artefacts.

The frozen repositories preserve original reference spans, held-out predictions, positive corpus flags, prompts, analysis scripts, and the released detector source. Row-level spans for the detector-informed revision were not archived. Exact replay remains unavailable because the original study did not record analysed repository revisions, the detector revision used for the runs, raw API responses, request identifiers and complete metadata, retry histories, complete provider metadata, or the full \num{17450}-file corpus frame including files without flags. The materials support inspection of surviving outputs and partial reruns of the released pipeline; they cannot reconstruct the original requests and source states exactly.
